\documentclass[12pt]{article}
\usepackage[utf8]{inputenc}
\usepackage[T1]{fontenc}
\usepackage{hyperref}
\title{Tingal Language Resources for K-12 Curriculum Development}
\author{Compiled by Claude Sonnet 5 (claude-sonnet-5) for Andr\'as Kornai}
\date{}
\begin{document}
\maketitle
\section{Introduction}
Tingal (ISO 639-3 tie) is a small Nuba language of South Kordofan, Sudan. The family label ``Bookkeeping'' in the task identification is evidently a data error. Sources link the Tingal to the Rashad/Tegali cluster, which is Kordofanian, a branch of Niger-Congo. I could not confirm how the Tingal language relates to the language called Tegali, so the identification is provisional.
\section{Primary Educational Resources}
None found.
\section{Linguistic and Reference Materials}
\begin{itemize}
\item Wikipedia on Tegali (the page excerpt contained only citations, so I could not check what it says): \url{https://en.wikipedia.org/wiki/Tegali_language}
\item Quint, ``Do you speak Kordofanian?'': \url{https://halshs.archives-ouvertes.fr/halshs-00171745}
\end{itemize}
\section{Religious and Community Resources}
Joshua Project has a profile of the Tingal people (also called Kajaja): \url{https://joshuaproject.net/people_groups/15496/SU}.
\section{Digital and Audio Resources}
None found.
\section{Curriculum Development Context}
Little exists. Many Tingal also speak Sudanese Arabic, which dominates trade. Any curriculum would need new documentation work first.
\end{document}